\part{Dieselben über die Zeit}

\chapter{Jeder sein eigener Vergleich}

\section{Was Paneldaten können}

\begin{merksatz}
\boxtitel{Dieselben Versicherten, sechs Jahre}
$5000$ Versicherte, jeder sechs Jahre lang beobachtet. Die Zuzahlung ändert sich über die Zeit --- durch Tarifwechsel und allgemeine Erhöhungen.

\begin{center}
{\small
\renewcommand{\arraystretch}{1.48}
\begin{tabular}{@{}P{3.6cm}P{2.6cm}@{}}
\toprule
\textbf{Verfahren} & \textbf{je Euro}\\
\midrule
gepoolt (alles in einen Topf) & $-0{,}302$\\
zwischen (Personenmittel) & $-0{,}418$\\
innerhalb (Abweichungen) & $-0{,}070$\\
innerhalb, mit Jahreseffekten & $\mathbf{-0{,}099}$\\
\midrule
wahrer Wert & $\mathbf{-0{,}100}$\\
\bottomrule
\end{tabular}}
\end{center}

\begin{center}
\begin{minipage}{0.9\textwidth}
\textbf{Dieselben Daten, vier Antworten --- und nur die vierte stimmt.}
\end{minipage}
\end{center}
\end{merksatz}

\begin{definition}
\boxtitelm{Definition 7.1 --- das Modell mit festen Effekten}
\[ y_{it} = \beta x_{it} + \alpha_i + \delta_t + u_{it}, \qquad i = 1,\ldots,n,\quad t = 1,\ldots,T \]
Dabei ist $\alpha_i$ ein \Begriff{fester Effekt} --- alles an Person $i$, was sich über die Zeit nicht ändert, beobachtet oder nicht --- und $\delta_t$ ein \Begriff{Zeiteffekt}, der alle Personen in Periode $t$ gleichermaßen trifft.
\end{definition}

\begin{merksatz}
\boxtitel{Satz 7.2 --- die Innerhalb-Transformation}
Zieht man von jeder Größe ihr Personenmittel ab,
\[ \tilde y_{it} = y_{it} - \bar y_i, \qquad \tilde x_{it} = x_{it} - \bar x_i , \]
so verschwindet $\alpha_i$, und die Regression von $\tilde y$ auf $\tilde x$ ist konsistent --- sofern $u_{it}$ mit $x_{is}$ für alle $s$ unkorreliert ist.

\schritt{Beweis}
\[ \bar y_i = \beta\bar x_i + \alpha_i + \bar\delta + \bar u_i \quad\Longrightarrow\quad \tilde y_{it} = \beta\tilde x_{it} + (\delta_t-\bar\delta) + \tilde u_{it} \]
Der feste Effekt kürzt sich heraus. Da $\tilde u_{it}$ eine Linearkombination der $u_{is}$ ist, braucht man deren Unkorreliertheit mit $x$ in allen Perioden. $\blacksquare$

\begin{center}
\begin{minipage}{0.9\textwidth}
\textbf{Man vergleicht jeden Versicherten nur mit sich selbst --- und alles, was ihn zeitkonstant von anderen unterscheidet, fällt heraus.}
\end{minipage}
\end{center}

\begin{center}
\begin{minipage}{0.88\textwidth}
\kursiv{\textbf{Im Leitfall ist die Gesundheitsneigung der feste Effekt.} Sie war in Kapitel~1 die unbeobachtete Störgröße, für die man weder kontrollieren noch ein Instrument finden konnte. Hier verschwindet sie ohne Messung und ohne Instrument --- nur deshalb, weil sie sich innerhalb der sechs Jahre nicht ändert. Diese eine Annahme ersetzt alle anderen.}
\end{minipage}
\end{center}
\end{merksatz}

\section{Die drei Regressionen}

\begin{merksatz}
\boxtitel{Woher die vier Antworten kommen}
\begin{center}
{\small
\renewcommand{\arraystretch}{1.55}
\begin{tabular}{@{}P{2.4cm}P{6.2cm}@{}}
\toprule
\textbf{Gepoolt} & ignoriert $\alpha_i$ --- die Verzerrung aus Kapitel~1.\\
\textbf{Zwischen} & regressiert Personenmittel auf Personenmittel; enthält \kursiv{nur} $\alpha_i$ und ist am stärksten verzerrt.\\
\textbf{Innerhalb} & beseitigt $\alpha_i$ --- aber nicht $\delta_t$: Zuzahlung und Besuche steigen beide über die Jahre, und der Trend wird der Zuzahlung zugeschrieben.\\
\textbf{Mit Jahreseffekten} & beseitigt beides.\\
\bottomrule
\end{tabular}}
\end{center}

\begin{center}
\begin{minipage}{0.88\textwidth}
\kursiv{\textbf{Die dritte Zeile ist eine Warnung, die man leicht übersieht.} Feste Effekte für Personen sind das bekannte Werkzeug; feste Effekte für Zeitpunkte sind ebenso nötig, sobald irgendetwas alle zugleich betrifft --- Preise, Gesetze, Grippewellen. Ein Modell mit beiden heißt zweiseitig, und es ist der Normalfall, nicht die Ausnahme.}
\end{minipage}
\end{center}

\schritt{Erste Differenzen}
Statt des Personenmittels kann man den Vorjahreswert abziehen:
\[ \Delta y_{it} = \beta\,\Delta x_{it} + \Delta\delta_t + \Delta u_{it} \]
\kursiv{Im Leitfall: $-0{,}0995$ --- bei $T = 2$ identisch mit der Innerhalb-Schätzung, sonst eine Alternative mit anderen Eigenschaften (Aufgabe 7.2).}
\end{merksatz}

\section{Der Bruchfall}

\begin{warnung}
\boxtitelr{Was feste Effekte kosten}
\begin{center}
{\small
\renewcommand{\arraystretch}{1.52}
\begin{tabular}{@{}P{2.8cm}P{5.8cm}@{}}
\toprule
\textbf{Zeitkonstantes} & Alles, was sich nicht ändert, ist nicht schätzbar --- Geschlecht, Geburtsjahr, Bildung. Es ist im festen Effekt enthalten.\\
\textbf{Wenig Bewegung} & Nur die Veränderung innerhalb der Person zählt. Im Leitfall sind das ein Drittel der Gesamtvarianz.\\
\textbf{Messfehler} & Er wird verstärkt: Der wahre Teil der Schwankung schrumpft, der Fehler bleibt.\\
\bottomrule
\end{tabular}}
\end{center}

\schritt{Zur dritten Zeile}
Die Zuzahlung sei mit einem Fehler der Streuung $2$ erfasst:

\begin{center}
{\small
\renewcommand{\arraystretch}{1.48}
\begin{tabular}{@{}P{2.6cm}P{2.2cm}P{2.6cm}@{}}
\toprule
\textbf{Verfahren} & \textbf{ohne Fehler} & \textbf{mit Fehler}\\
\midrule
gepoolt & $-0{,}302$ & $-0{,}282$\\
innerhalb & $-0{,}099$ & $\mathbf{-0{,}059}$\\
\bottomrule
\end{tabular}}
\end{center}

\begin{center}
\begin{minipage}{0.9\textwidth}
\textbf{Der Messfehler dämpft die Innerhalb-Schätzung um vierzig Prozent, die gepoolte nur um sieben.}
\end{minipage}
\end{center}

\begin{center}
\begin{minipage}{0.88\textwidth}
\kursiv{Der Grund ist die Dämpfungsformel aus Kapitel~1: Der Faktor ist $\sigma_x^{2}/(\sigma_x^{2}+\sigma_e^{2})$, und die Innerhalb-Transformation verkleinert $\sigma_x^{2}$ auf ein Drittel, während $\sigma_e^{2}$ bleibt. \textbf{Wer feste Effekte benutzt, bezahlt mit Empfindlichkeit gegen Messfehler} --- und wer dann einen kleineren Effekt findet, kann nicht ohne weiteres sagen, ob die Störgröße oder der Messfehler verschwunden ist.}
\end{minipage}
\end{center}
\end{warnung}

\section{Übungsaufgaben}

\begin{aufgabe}
\boxtitel{Aufgabe 7.1 --- Zwei Perioden von Hand}
Ein Versicherter zahlt im ersten Jahr $10$ Euro und geht $7$-mal zum Arzt, im zweiten Jahr $20$ Euro und $6$-mal. Ein zweiter --- kränker und in einem teureren Tarif --- zahlt $15$ und $35$ Euro und geht $9$- und $7$-mal. Schätzen Sie $\beta$ mit festen Effekten.

\loesung
\schritt{Schritt 1 --- Die Differenzen}
\begin{center}
{\small
\renewcommand{\arraystretch}{1.48}
\begin{tabular}{@{}P{1.6cm}P{2.0cm}P{2.0cm}P{2.4cm}@{}}
\toprule
\textbf{Person} & $\Delta x$ & $\Delta y$ & $\Delta y/\Delta x$\\
\midrule
$1$ & $10$ & $-1$ & $-0{,}100$\\
$2$ & $20$ & $-2$ & $-0{,}100$\\
\bottomrule
\end{tabular}}
\end{center}

\schritt{Schritt 2 --- Die Regression durch den Ursprung}
\[ \hat\beta = \frac{\sum\Delta x\,\Delta y}{\sum(\Delta x)^{2}} = \frac{-10-40}{100+400} = \mathbf{-0{,}100} \]

\schritt{Schritt 3 --- Was der gepoolte Schätzer sagt}
Vier Punkte $(10,7)$, $(20,6)$, $(15,9)$, $(35,7)$ mit $\bar x = 20$, $\bar y = 7{,}25$:
\[ \sum(x-\bar x)(y-\bar y) = (-10)(-0{,}25)+0+(-5)(1{,}75)+15(-0{,}25) = -10 \]
\[ \sum(x-\bar x)^{2} = 100+0+25+225 = 350 \quad\Longrightarrow\quad \hat\beta_{\text{gepoolt}} = \frac{-10}{350} = -0{,}029 \]

\begin{center}
\begin{minipage}{0.9\textwidth}
\textbf{Gepoolt: fast nichts. Mit festen Effekten: $-0{,}100$. Innerhalb jeder Person ist die Wirkung eindeutig --- zwischen ihnen verdeckt sie das Niveau.}
\end{minipage}
\end{center}

\begin{center}
\begin{minipage}{0.88\textwidth}
\kursiv{\textbf{Person~2 geht trotz höherer Zuzahlung öfter}, weil sie kränker ist. Der gepoolte Schätzer sieht das als Beleg, dass Zuzahlung kaum wirkt; die Innerhalb-Schätzung vergleicht nur die beiden Jahre jeder Person und findet in beiden dieselbe Steigung. Hier ist die Verzerrung nach oben --- das Gegenteil von Kapitel~1, weil der feste Effekt diesmal mit der Zuzahlung \kursiv{positiv} zusammenhängt. Die Richtung ist nicht vorhersagbar, nur die Tatsache.}
\end{minipage}
\end{center}
\end{aufgabe}

\begin{aufgabe}
\boxtitel{Aufgabe 7.2 --- Innerhalb oder Differenzen?}
Bei $T > 2$ liefern Innerhalb-Transformation und erste Differenzen verschiedene Schätzer. Wann welcher?

\loesung
\schritt{Schritt 1 --- Die Annahme}
Beide brauchen, dass $u_{it}$ mit $x_{is}$ unkorreliert ist --- aber in verschiedener Stärke:
\begin{center}
{\small
\renewcommand{\arraystretch}{1.5}
\begin{tabular}{@{}P{2.4cm}P{6.2cm}@{}}
\toprule
\textbf{Innerhalb} & für \kursiv{alle} $s$ --- strenge Exogenität.\\
\textbf{Differenzen} & nur für $s = t$ und $s = t-1$.\\
\bottomrule
\end{tabular}}
\end{center}

\schritt{Schritt 2 --- Die Effizienz}
\begin{center}
{\small
\renewcommand{\arraystretch}{1.5}
\begin{tabular}{@{}P{2.4cm}P{6.2cm}@{}}
\toprule
\textbf{Innerhalb} & effizient, wenn $u_{it}$ unkorreliert über die Zeit ist.\\
\textbf{Differenzen} & effizient, wenn $u_{it}$ eine Irrfahrt ist --- die Differenzen sind dann unkorreliert.\\
\bottomrule
\end{tabular}}
\end{center}

\schritt{Schritt 3 --- Die Regel}
\begin{center}
\begin{minipage}{0.9\textwidth}
\textbf{Differenzen, wenn die Störung Gedächtnis hat; Innerhalb, wenn nicht. Weichen beide deutlich ab, ist die Annahme der strengen Exogenität verletzt.}
\end{minipage}
\end{center}

\begin{center}
\begin{minipage}{0.88\textwidth}
\kursiv{\textbf{Die letzte Aussage ist ein Test.} Strenge Exogenität verbietet, dass ein heutiger Schock die morgige Zuzahlung ändert --- im Leitfall: dass ein Krankheitsjahr zum Tarifwechsel führt. Das ist eine starke Annahme, und ihre Verletzung heißt Rückkopplung. Kapitel~8 behandelt Modelle, die sie zulassen, und bezahlt dafür mit Instrumenten.}
\end{minipage}
\end{center}
\end{aufgabe}

\begin{aufgabe}
\boxtitel{Aufgabe 7.3 --- Zufällige Effekte}
Wann darf man $\alpha_i$ als Zufallsgröße behandeln statt als festen Parameter --- und was gewinnt man?

\loesung
\schritt{Schritt 1 --- Das Modell}
Man nimmt an, $\alpha_i$ sei zufällig mit Varianz $\sigma_\alpha^{2}$ und \kursiv{unkorreliert mit $x$}. Dann hat die Störung $\alpha_i + u_{it}$ die Korrelationsstruktur
\[ \operatorname{Corr}(\alpha_i+u_{it},\ \alpha_i+u_{is}) = \frac{\sigma_\alpha^{2}}{\sigma_\alpha^{2}+\sigma_u^{2}} \quad\text{für } s\ne t , \]
und man schätzt mit verallgemeinerten kleinsten Quadraten.

\schritt{Schritt 2 --- Der Gewinn}
\begin{center}
{\small
\renewcommand{\arraystretch}{1.5}
\begin{tabular}{@{}P{2.8cm}P{5.8cm}@{}}
\toprule
\textbf{Zeitkonstantes} & wird schätzbar.\\
\textbf{Effizienz} & nutzt Innerhalb- \kursiv{und} Zwischen-Variation.\\
\bottomrule
\end{tabular}}
\end{center}

\schritt{Schritt 3 --- Der Preis}
\begin{center}
\begin{minipage}{0.9\textwidth}
\textbf{Die Annahme, dass $\alpha_i$ mit $x$ unkorreliert ist, ist im Leitfall genau die, die Kapitel~1 widerlegt hat.}
\end{minipage}
\end{center}

\schritt{Schritt 4 --- Der Test}
\begin{center}
\begin{minipage}{0.88\textwidth}
\kursiv{Man vergleiche die Schätzungen mit festen und mit zufälligen Effekten. \textbf{Stimmen sie überein, so ist die Unkorreliertheit haltbar und das sparsamere Modell erlaubt}; weichen sie ab, so sind die zufälligen Effekte inkonsistent. Das ist der Hausman-Test --- und im Leitfall fällt er eindeutig aus: $-0{,}302$ gegen $-0{,}099$. In der Praxis ist das der häufigste Ausgang, und die festen Effekte gewinnen.}
\end{minipage}
\end{center}
\end{aufgabe}

\begin{geschichte}
\boxtitelg{Von Fisher zu den Panels}
Die Innerhalb-Transformation ist die Blockbildung aus Fishers Versuchsplanung der 1920er Jahre: Jeder Block ist eine Person, und verglichen wird nur innerhalb der Blöcke.

\begin{center}
\begin{minipage}{0.9\textwidth}
\kursiv{\Person{Yair Mundlak} zeigte 1961 an landwirtschaftlichen Betrieben, dass die gepoolte Schätzung der Produktionsfunktion durch unbeobachtete Managementqualität verzerrt ist --- und dass feste Effekte das beheben. \textbf{Es war die erste ökonometrische Anwendung, und die Managementqualität ist bis heute das Lehrbuchbeispiel für einen festen Effekt.}}
\end{minipage}
\end{center}

\schritt{Die großen Panels}
\begin{center}
\begin{minipage}{0.9\textwidth}
\kursiv{Das amerikanische Einkommenspanel begann 1968, das deutsche Sozio-oekonomische Panel 1984 --- \textbf{beide verfolgen Haushalte über inzwischen mehr als vierzig Jahre} und sind die Grundlage fast aller Aussagen über Einkommensmobilität, Erwerbsverläufe und Gesundheit im Lebenslauf. Ihr Wert liegt nicht in der Zahl der Beobachtungen, sondern darin, dass es dieselben sind.}
\end{minipage}
\end{center}

\schritt{Und die Gegenseite}
\begin{center}
\begin{minipage}{0.9\textwidth}
\kursiv{\Person{Zvi Griliches} und \Person{Jerry Hausman} wiesen 1986 nach, was der Bruchfall zeigt: \textbf{Feste Effekte verstärken Messfehler}, und viele der kleineren Schätzungen, die sie liefern, sind nicht bereinigt, sondern gedämpft. Der Aufsatz trägt den Titel \kursiv{Errors in Variables in Panel Data} und ist die nüchterne Ergänzung zu Mundlaks Begeisterung.}
\end{minipage}
\end{center}
\end{geschichte}

\begin{forschung}
\boxtitelp{Panels heute}
\begin{center}
{\small
\renewcommand{\arraystretch}{1.45}
\begin{tabular}{@{}P{3.2cm}P{5.0cm}@{}}
\toprule
\textbf{Gestaffelte Einführung} & Was feste Effekte schätzen, wenn eine Maßnahme zu verschiedenen Zeiten eingeführt wird.\\
\textbf{Interaktive Effekte} & Feste Effekte, die mit der Zeit wechselwirken --- Faktormodelle.\\
\textbf{Verwaltungsdaten} & Panels aus Registern mit Millionen Personen über Jahrzehnte.\\
\bottomrule
\end{tabular}}
\end{center}

\schritt{Zur ersten Zeile}
\begin{center}
\begin{minipage}{0.88\textwidth}
\kursiv{Seit 2018 ist bekannt, dass das zweiseitige Modell bei gestaffelter Einführung und heterogenen Wirkungen etwas Unsinniges schätzt --- ein gewogenes Mittel, in dem manche Gewichte \kursiv{negativ} sind. \textbf{Hunderte veröffentlichter Studien beruhen auf dieser Spezifikation}, und die Literatur zur Berichtigung ist inzwischen selbst unübersichtlich. Kapitel~9 behandelt die doppelte Differenz, auf die das alles zurückgeht.}
\end{minipage}
\end{center}
\end{forschung}

\begin{merksatz}
\boxtitel{Was Sie jetzt können}
\begin{center}
\renewcommand{\arraystretch}{1.5}
\begin{tabular}{@{}P{3.0cm}P{5.6cm}@{}}
\toprule
\textbf{Definition 7.1} & Personen- und Zeiteffekte\\
\textbf{Satz 7.2} & Innerhalb-Transformation beseitigt $\alpha_i$\\
\textbf{Vier Antworten} & gepoolt, zwischen, innerhalb, zweiseitig\\
\textbf{Der Bruchfall} & Messfehler wird verstärkt\\
\textbf{Hausman} & fest gegen zufällig\\
\bottomrule
\end{tabular}
\end{center}

\begin{center}
\begin{minipage}{0.88\textwidth}
\kursiv{Feste Effekte verlangen, dass heutige Schocks die morgige Zuzahlung nicht beeinflussen. Das nächste Kapitel lässt das zu --- und findet, dass dann selbst die Innerhalb-Schätzung verzerrt ist, und zwar in berechenbarer Richtung.}
\end{minipage}
\end{center}
\end{merksatz}
